% ============================================================
\chapter{March\'es Financiers --- Mod\'elisation}
\label{ch:marches}
% ============================================================

En 1900, Louis Bachelier soutient sa th\`ese \`a la Sorbonne sous la
direction d'Henri Poincar\'e. Son sujet~: la \og th\'eorie de la
sp\'eculation\fg, o\`u il propose de mod\'eliser les cours boursiers par un
processus stochastique. Poincar\'e, pourtant visionnaire, juge le travail
\og pas sans int\'er\^et\fg{} mais lui accorde une mention honorable plut\^ot
que tr\`es honorable. Il faudra soixante ans, l'arriv\'ee de Samuelson,
Black, Scholes et Merton, pour que les id\'ees de Bachelier deviennent la
pierre angulaire de la finance moderne. Aujourd'hui, les march\'es
financiers brassent des milliers de milliards par jour, et les mod\`eles
math\'ematiques qui les d\'ecrivent sont au c\oe ur de l'\'economie mondiale.

% ------------------------------------------------------------
\section{Introduction aux march\'es financiers}
% ------------------------------------------------------------

Les march\'es financiers constituent le lieu d'\'echange d'actifs dont la
valeur est incertaine. Comprendre leur fonctionnement et les mod\'eliser
math\'ematiquement est le point de d\'epart de toute la finance quantitative.

\begin{definition}[March\'e financier]
Un \textbf{march\'e financier} est un syst\`eme organis\'e permettant
l'\'echange d'actifs financiers entre agents \'economiques. On distingue~:
\begin{itemize}
    \item Le \textbf{march\'e primaire}~: \'emission de nouveaux titres.
    \item Le \textbf{march\'e secondaire}~: \'echange de titres d\'ej\`a \'emis.
\end{itemize}
\end{definition}

\subsection{Classes d'actifs}

\begin{definition}[Actif financier]
Un \textbf{actif financier} est un contrat donnant droit \`a des flux
de tr\'esorerie futurs. Les principales classes sont~:
\begin{enumerate}
    \item \textbf{Actions} (equities)~: parts de propri\'et\'e d'une entreprise.
    \item \textbf{Obligations} (bonds)~: titres de cr\'eance \`a revenu fixe.
    \item \textbf{Produits d\'eriv\'es} (derivatives)~: contrats dont la valeur
          d\'epend d'un actif sous-jacent.
    \item \textbf{Devises} (currencies)~: march\'e des changes (forex).
    \item \textbf{Mati\`eres premi\`eres} (commodities)~: or, p\'etrole, bl\'e.
\end{enumerate}
\end{definition}

\subsection{Produits d\'eriv\'es fondamentaux}

\begin{definition}[Option europ\'eenne]
Une \textbf{option d'achat (call) europ\'eenne} de prix d'exercice $K$ et
de maturit\'e $T$ sur un sous-jacent $S$ donne \`a son d\'etenteur le droit
(mais non l'obligation) d'acheter $S$ au prix $K$ \`a la date $T$. Son
payoff est~:
\[
    (S_T - K)^+ = \max(S_T - K, 0).
\]
De m\^eme, une \textbf{option de vente (put) europ\'eenne} a pour payoff~:
\[
    (K - S_T)^+ = \max(K - S_T, 0).
\]
\end{definition}

\begin{definition}[Contrat forward]
Un \textbf{contrat forward} (ou contrat \`a terme) est un accord d'acheter
ou de vendre un actif \`a une date future $T$ \`a un prix $F$ fix\'e
aujourd'hui. Son payoff pour l'acheteur est $S_T - F$.
\end{definition}

\section{Mod\'elisation en temps discret}

\subsection{Mod\`ele \`a une p\'eriode}

Consid\'erons un march\'e \`a deux dates $t=0$ et $t=1$, avec~:
\begin{itemize}
    \item Un actif sans risque de prix $B_0 = 1$, $B_1 = 1+r$.
    \item Un actif risqu\'e de prix $S_0$ \`a $t=0$, pouvant prendre les
          valeurs $S_1^u = uS_0$ (hausse) ou $S_1^d = dS_0$ (baisse)
          avec $0 < d < 1+r < u$.
\end{itemize}

\begin{definition}[Strat\'egie de trading]
Une \textbf{strat\'egie de trading} est un couple $(\phi^0, \phi^1) \in \R^2$
o\`u $\phi^0$ repr\'esente la quantit\'e d'actif sans risque et $\phi^1$
la quantit\'e d'actif risqu\'e d\'etenues. La valeur du portefeuille est~:
\[
    V_t = \phi^0 B_t + \phi^1 S_t, \quad t \in \{0,1\}.
\]
\end{definition}

\begin{definition}[Opportunit\'e d'arbitrage]
Une \textbf{opportunit\'e d'arbitrage} est une strat\'egie $(\phi^0, \phi^1)$
telle que~:
\[
    V_0 \leq 0, \quad V_1 \geq 0 \text{ p.s.}, \quad \PP(V_1 > 0) > 0.
\]
Autrement dit, un gain possible sans risque de perte et sans mise de fonds
initiale.
\end{definition}

\begin{theoreme}[Absence d'arbitrage --- mod\`ele binomial \`a une p\'eriode]
Le mod\`ele binomial \`a une p\'eriode est sans arbitrage si et seulement si
$d < 1+r < u$.
\end{theoreme}

\begin{proof}
Supposons qu'il existe un arbitrage $(\phi^0, \phi^1)$ avec $V_0 = 0$. Alors~:
\[
    \phi^0 + \phi^1 S_0 = 0 \implies \phi^0 = -\phi^1 S_0.
\]
Les valeurs terminales sont~:
\begin{align*}
    V_1^u &= \phi^1(uS_0 - (1+r)S_0) = \phi^1 S_0(u - 1 - r), \\
    V_1^d &= \phi^1(dS_0 - (1+r)S_0) = \phi^1 S_0(d - 1 - r).
\end{align*}
Pour que $V_1^u \geq 0$ et $V_1^d \geq 0$ simultan\'ement avec au moins
une in\'egalit\'e stricte, il faut que $u - 1 - r$ et $d - 1 - r$ soient
de m\^eme signe non nul, ce qui contredit $d < 1+r < u$.
\end{proof}

\subsection{Probabilit\'e risque-neutre}

\begin{theoreme}[Mesure risque-neutre]
\label{thm:risque_neutre}
Sous la condition AOA ($d < 1+r < u$), il existe une unique mesure de
probabilit\'e $\mathbb{Q}$ sur $\{\omega_u, \omega_d\}$ telle que le
prix actualis\'e de l'actif risqu\'e soit une $\mathbb{Q}$-martingale~:
\[
    S_0 = \frac{1}{1+r}\E^\mathbb{Q}[S_1].
\]
La probabilit\'e risque-neutre est donn\'ee par~:
\[
    q = \frac{(1+r) - d}{u - d}, \quad 1-q = \frac{u - (1+r)}{u - d}.
\]
\end{theoreme}

\begin{formules}[title=\textbf{Prix risque-neutre d'un actif contingent}]
Le prix \`a $t=0$ d'un actif contingent de payoff $H = h(S_1)$ est~:
\[
    V_0 = \frac{1}{1+r}\E^\mathbb{Q}[H]
        = \frac{1}{1+r}\bigl[q \cdot h(uS_0) + (1-q) \cdot h(dS_0)\bigr].
\]
\end{formules}

\begin{exemple}
Soit $S_0 = 100$, $u = 1.2$, $d = 0.9$, $r = 0.05$, $K = 100$.

La probabilit\'e risque-neutre est~:
\[
    q = \frac{1.05 - 0.9}{1.2 - 0.9} = \frac{0.15}{0.30} = 0.5.
\]
Le prix du call est~:
\[
    C_0 = \frac{1}{1.05}\bigl[0.5 \times (120-100)^+ + 0.5 \times (90-100)^+\bigr]
        = \frac{0.5 \times 20}{1.05} \approx 9.52.
\]
\end{exemple}

\section{Mod\`ele binomial multi-p\'eriodes (CRR)}

\begin{definition}[Mod\`ele de Cox--Ross--Rubinstein]
Le \textbf{mod\`ele CRR} \`a $N$ p\'eriodes est d\'efini par~:
\begin{itemize}
    \item $S_{n+1} = S_n \cdot Y_{n+1}$ o\`u $Y_{n+1} \in \{u, d\}$
          sont i.i.d. sous $\mathbb{Q}$.
    \item Le prix actualis\'e $\tilde{S}_n = S_n / (1+r)^n$ est une
          $\mathbb{Q}$-martingale.
    \item Apr\`es $n$ p\'eriodes~: $S_n = S_0 u^j d^{n-j}$ pour
          $j$ hausses parmi $n$.
\end{itemize}
\end{definition}

\begin{theoreme}[Formule de pricing CRR]
Le prix d'un call europ\'een dans le mod\`ele CRR est~:
\[
    C_0 = \frac{1}{(1+r)^N} \sum_{j=0}^{N} \binom{N}{j} q^j (1-q)^{N-j}
    \bigl(S_0 u^j d^{N-j} - K\bigr)^+.
\]
\end{theoreme}

\begin{intuition}[title=\textbf{Convergence vers Black--Scholes}]
Lorsque $N \to \infty$ avec le param\'etrage $u = e^{\sigma\sqrt{\Delta t}}$,
$d = e^{-\sigma\sqrt{\Delta t}}$, $\Delta t = T/N$, la formule CRR converge
vers la formule de Black--Scholes (chapitre 2). C'est une cons\'equence du
th\'eor\`eme central limite.
\end{intuition}

\subsection{Calibration des param\`etres CRR}

Pour que le mod\`ele discret approche le mod\`ele continu avec volatilit\'e
$\sigma$ et taux $r$, on pose~:
\begin{align}
    u &= e^{\sigma\sqrt{\Delta t}}, \quad d = e^{-\sigma\sqrt{\Delta t}} = 1/u, \\
    q &= \frac{e^{r\Delta t} - d}{u - d}.
\end{align}

\section{Th\'eor\`emes fondamentaux de l'\'evaluation d'actifs}

\begin{theoreme}[Premier th\'eor\`eme fondamental (FTAP)]
\label{thm:ftap1}
Un march\'e est sans arbitrage si et seulement s'il existe une mesure
de probabilit\'e $\mathbb{Q}$ \'equivalente \`a $\PP$ sous laquelle les
prix actualis\'es des actifs sont des martingales.
\end{theoreme}

\begin{theoreme}[Second th\'eor\`eme fondamental]
\label{thm:ftap2}
Un march\'e sans arbitrage est \textbf{complet} (tout actif contingent est
r\'eplicable) si et seulement si la mesure martingale $\mathbb{Q}$ est unique.
\end{theoreme}

\begin{remarque}
Dans le mod\`ele binomial, le march\'e est complet car il y a deux \'etats
de la nature et deux actifs. En revanche, dans un mod\`ele trinomial
(trois \'etats), le march\'e est incomplet et il existe une infinit\'e
de mesures martingales.
\end{remarque}

\section{R\'eplication et couverture}

\begin{proposition}[Portefeuille de r\'eplication --- une p\'eriode]
Dans le mod\`ele binomial \`a une p\'eriode, l'actif contingent
$H = h(S_1)$ est r\'epliqu\'e par $(\phi^0, \phi^1)$ avec~:
\[
    \phi^1 = \frac{h(uS_0) - h(dS_0)}{(u-d)S_0}, \quad
    \phi^0 = \frac{1}{1+r}\bigl[h(uS_0) - \phi^1 uS_0\bigr].
\]
La quantit\'e $\phi^1$ est appel\'ee le \textbf{delta} de l'option.
\end{proposition}

\section{Impl\'ementation Python}

\begin{codeblock}[title=\textbf{Mod\`ele binomial CRR en Python}]
\begin{minted}{python}
import numpy as np

def crr_european(S0, K, T, r, sigma, N, option='call'):
    """Prix d'une option europeenne par le modele CRR."""
    dt = T / N
    u = np.exp(sigma * np.sqrt(dt))
    d = 1 / u
    q = (np.exp(r * dt) - d) / (u - d)
    discount = np.exp(-r * dt)

    # Prix terminaux
    j = np.arange(N + 1)
    ST = S0 * u**j * d**(N - j)

    if option == 'call':
        payoff = np.maximum(ST - K, 0)
    else:
        payoff = np.maximum(K - ST, 0)

    # Backward induction
    for n in range(N - 1, -1, -1):
        payoff = discount * (q * payoff[1:] + (1 - q) * payoff[:-1])

    return payoff[0]

# Exemple
S0, K, T, r, sigma = 100, 100, 1.0, 0.05, 0.20
for N in [10, 50, 100, 500, 1000]:
    price = crr_european(S0, K, T, r, sigma, N)
    print(f"N={N:4d}: C = {price:.4f}")
\end{minted}
\end{codeblock}

\begin{codeblock}[title=\textbf{Parit\'e call-put --- v\'erification}]
\begin{minted}{python}
def verify_put_call_parity(S0, K, T, r, sigma, N):
    """Verifie la parite call-put: C - P = S0 - K*exp(-rT)."""
    C = crr_european(S0, K, T, r, sigma, N, 'call')
    P = crr_european(S0, K, T, r, sigma, N, 'put')
    parity = S0 - K * np.exp(-r * T)
    print(f"C - P = {C - P:.6f}")
    print(f"S0 - K*exp(-rT) = {parity:.6f}")
    print(f"Erreur = {abs(C - P - parity):.2e}")

verify_put_call_parity(100, 100, 1.0, 0.05, 0.20, 500)
\end{minted}
\end{codeblock}

\section{Parit\'e call-put}

\begin{theoreme}[Parit\'e call-put]
Pour des options europ\'eennes de m\^eme strike $K$ et maturit\'e $T$~:
\[
    C - P = S_0 - K e^{-rT}.
\]
\end{theoreme}

\begin{proof}
Consid\'erons le portefeuille $\Pi = C - P - S_0 + Ke^{-rT}$. \`A maturit\'e~:
\[
    \Pi_T = (S_T - K)^+ - (K - S_T)^+ - S_T + K = S_T - K - S_T + K = 0.
\]
Par AOA, $\Pi_0 = 0$, d'o\`u le r\'esultat.
\end{proof}

\section{Hypoth\`eses de march\'e}

\begin{attention}[title=\textbf{Hypoth\`eses simplificatrices}]
Les mod\`eles classiques reposent sur des hypoth\`eses fortes~:
\begin{enumerate}
    \item March\'e sans friction (pas de co\^uts de transaction).
    \item Trading continu possible.
    \item Vente \`a d\'ecouvert autoris\'ee sans restriction.
    \item Taux d'emprunt \'egal au taux de pr\^et.
    \item Pas de risque de contrepartie.
\end{enumerate}
La relaxation de ces hypoth\`eses conduit \`a des mod\`eles plus r\'ealistes
mais significativement plus complexes.
\end{attention}

\section{Faits stylis\'es des rendements financiers}

Les rendements logarithmiques $r_t = \ln(S_t/S_{t-1})$ pr\'esentent
des propri\'et\'es empiriques bien document\'ees~:

\begin{enumerate}
    \item \textbf{Queues \'epaisses}~: la distribution des rendements a des
          queues plus lourdes que la loi normale (kurtosis exc\'edentaire).
    \item \textbf{Asym\'etrie}~: l\'eg\`ere asym\'etrie n\'egative.
    \item \textbf{Agr\'egation gaussienne}~: les rendements sur longue
          p\'eriode se rapprochent de la loi normale.
    \item \textbf{Absence d'autocorr\'elation}~: les rendements sont
          quasi non corr\'el\'es.
    \item \textbf{Clustering de volatilit\'e}~: les p\'eriodes de forte
          volatilit\'e sont regroup\'ees.
    \item \textbf{Effet de levier}~: corr\'elation n\'egative entre
          rendements et volatilit\'e.
\end{enumerate}

\begin{codeblock}[title=\textbf{Analyse des faits stylis\'es}]
\begin{minted}{python}
import numpy as np
from scipy import stats

def stylized_facts(returns):
    """Analyse des faits stylises des rendements."""
    print(f"Moyenne:    {np.mean(returns):.6f}")
    print(f"Ecart-type: {np.std(returns):.6f}")
    print(f"Skewness:   {stats.skew(returns):.4f}")
    print(f"Kurtosis:   {stats.kurtosis(returns):.4f}")

    # Test de normalite (Jarque-Bera)
    jb_stat, jb_pval = stats.jarque_bera(returns)
    print(f"Jarque-Bera: stat={jb_stat:.2f}, p={jb_pval:.4e}")

    # Autocorrelation des rendements et des carres
    n = len(returns)
    r_centered = returns - np.mean(returns)
    acf1 = np.correlate(r_centered[:-1], r_centered[1:]) / (n * np.var(returns))
    r2 = r_centered**2
    acf1_sq = np.correlate(r2[:-1], r2[1:]) / (n * np.var(r2))
    print(f"ACF(1) rendements: {acf1[0]:.4f}")
    print(f"ACF(1) carres:     {acf1_sq[0]:.4f}")

# Simulation GBM pour comparaison
np.random.seed(42)
S0, mu, sigma, dt, N = 100, 0.08, 0.20, 1/252, 2520
Z = np.random.standard_normal(N)
log_returns = (mu - 0.5*sigma**2)*dt + sigma*np.sqrt(dt)*Z
stylized_facts(log_returns)
\end{minted}
\end{codeblock}

\section{Exercices}

\begin{exercice}[Mod\`ele binomial \`a une p\'eriode]
Soit $S_0 = 50$, $u = 1.3$, $d = 0.8$, $r = 0.04$.
\begin{enumerate}
    \item V\'erifier la condition AOA.
    \item Calculer la probabilit\'e risque-neutre $q$.
    \item Calculer le prix d'un call de strike $K = 55$.
    \item D\'eterminer le portefeuille de r\'eplication $(\phi^0, \phi^1)$.
    \item V\'erifier que le prix du put satisfait la parit\'e call-put.
\end{enumerate}
\end{exercice}

\begin{exercice}[Convergence CRR]
\begin{enumerate}
    \item Impl\'ementer le mod\`ele CRR et tracer le prix du call en
          fonction de $N$ pour $N = 1, 2, \ldots, 200$.
    \item Observer la convergence et les oscillations pair/impair.
    \item Comparer avec la formule de Black--Scholes (chapitre 2).
\end{enumerate}
\end{exercice}

\begin{exercice}[March\'e incomplet]
Consid\'erons un mod\`ele trinomial~: $S_1 \in \{uS_0, mS_0, dS_0\}$
avec $d < m < u$.
\begin{enumerate}
    \item Montrer que le march\'e est incomplet.
    \item Caract\'eriser l'ensemble des mesures martingales.
    \item En d\'eduire les bornes de non-arbitrage pour le prix d'un call.
\end{enumerate}
\end{exercice}

\begin{exercice}[Options am\'ericaines dans le mod\`ele CRR]
\begin{enumerate}
    \item Modifier l'algorithme CRR pour pricer un put am\'ericain en
          int\'egrant la possibilit\'e d'exercice anticip\'e.
    \item Comparer les prix du put am\'ericain et europ\'een. V\'erifier
          que le put am\'ericain est toujours au moins aussi cher.
    \item D\'eterminer la fronti\`ere d'exercice optimal.
\end{enumerate}
\end{exercice}
